% ============================================================
%  Grayscale Image Colorization with GANs — Bachelor's Thesis
% ============================================================
\documentclass[12pt, a4paper, oneside]{report}

% ---- Encoding & language ----
\usepackage[utf8]{inputenc}
\usepackage[T1]{fontenc}
\usepackage[english]{babel}

% ---- Page layout ----
\usepackage[
    top=2.5cm, bottom=2.5cm,
    left=3cm,  right=2.5cm
]{geometry}
\usepackage{setspace}
\onehalfspacing

% ---- Math ----
\usepackage{amsmath}
\usepackage{amssymb}

% ---- Figures & tables ----
\usepackage{graphicx}
\usepackage{subcaption}   % side-by-side figures (grayscale | colorized | ground truth)
\usepackage{booktabs}     % clean tables for ablation results
\usepackage{float}

% ---- Code listings ----
\usepackage{listings}
\usepackage{xcolor}

% ---- References ----
\usepackage[backend=biber, style=ieee]{biblatex}
\addbibresource{references.bib}

% ---- Hyperlinks (keep last) ----
\usepackage[hidelinks]{hyperref}


% ============================================================
\begin{document}
% ============================================================


% ---- Title page ----
\begin{titlepage}
    \centering
    \vspace*{2cm}

    {\Large \textbf{Craiova University}}\\[0.5cm]
    {\large Faculty of Automation, Computers and Electronics}\\[3cm]

    {\Huge \textbf{Grayscale Image Colorization\\using Generative Adversarial Networks}}\\[2cm]

    {\large Bachelor's Thesis}\\[2cm]

    \begin{minipage}{0.45\textwidth}
        \begin{flushleft}
            \textbf{Supervisor:}\\
            Costin Badica
        \end{flushleft}
    \end{minipage}
    \begin{minipage}{0.45\textwidth}
        \begin{flushright}
            \textbf{Author:}\\
            Smarandache Andra-Maria
        \end{flushright}
    \end{minipage}

    \vfill
    {\large Craiova, 2026}
\end{titlepage}


% ---- Abstract ----
\begin{abstract}
% TODO: write last, after all results are in.
% Must state: (1) the problem, (2) your approach, (3) your novel contributions,
% (4) at least one quantitative result.
%
% Example structure:
% Automatic colorization of grayscale images is a challenging ill-posed problem...
% We propose a GAN-based system using a ResNet-18 U-Net generator and a PatchGAN
% discriminator operating in the LAB color space. We introduce two novel loss
% functions: Spatial Color Affinity Loss, which... and No Grey Loss, which...
% Experiments on a 13,000-image COCO subset show that our combined loss model
% achieves a PSNR of X dB and FID of Y, outperforming the baseline by Z%.
\end{abstract}

\tableofcontents
\listoffigures
\listoftables


% ============================================================
\chapter{Introduction}
% ============================================================

% What is the problem and why does it matter?
% What is your approach at a high level?
% What are your specific contributions?
% How is the thesis structured?

\section{Motivation}
\section{Problem Statement}
\section{Contributions}

% List your novel contributions explicitly — examiners look for this:
% 1. Spatial Color Affinity Loss
% 2. No Grey Loss
% 3. Adversarial robustness evaluation via FGSM/PGD

\section{Thesis Structure}


% ============================================================
\chapter{Related Work}
% ============================================================

% You must cite and discuss the foundational papers. Examiners check this.

\section{Early Approaches to Image Colorization}
% Discuss hand-crafted and scribble-based methods before deep learning.

\section{Deep Learning for Colorization}
% Zhang et al. (2016) — Colorful Image Colorization — the foundational paper.
% Larsson et al. (2016) — Learning Representations for Automatic Colorization.
% Iizuka et al. (2016) — Let there be Color.

\section{Generative Adversarial Networks}
% Goodfellow et al. (2014) — original GAN paper.
% Isola et al. (2017) — pix2pix, conditional GANs for image-to-image translation.
% The PatchGAN discriminator.

\section{Loss Functions in Image Synthesis}
% Perceptual loss (Johnson et al. 2016), Total Variation, L1 vs L2.
% Position your novel losses in relation to existing work.


% ============================================================
\chapter{Methodology}
% ============================================================

\section{Color Space}
% Why LAB instead of RGB.
% L = lightness (input), ab = color (output to predict).
% Normalization: L/50 - 1, ab/110.

\section{Dataset Preparation}
% Describe the COCO dataset and why we chose it.
% The grayscale image problem — automated detection + manual pass.
% The 16,000 subset strategy and why we needed that margin.
% Final split: 13,000 images, 80/20 train/validation.

\section{Model Architecture}

    \subsection{Generator: ResNet-18 U-Net}
    % Why ResNet-18 as backbone (pretrained on ImageNet).
    % U-Net skip connections — preserving spatial detail.
    % Input: 1 channel (L), output: 2 channels (ab).
    % First conv layer adaptation for grayscale input.

    \subsection{Discriminator: PatchGAN}
    % Why patch-level discrimination instead of image-level.
    % Encourages locally consistent texture and sharp edges.
    % Architecture: strided convolutions, BCEWithLogitsLoss on logits.

\section{Loss Functions}

    \subsection{Baseline Losses}
    % L1 loss (pixel fidelity, weight 100).
    % GAN loss (adversarial realism).
    % Total Variation loss (smoothness, artifact reduction).
    % Contrast loss.

    \subsection{Spatial Color Affinity Loss (Novel)}
    % The core idea: pixels with similar luminance should have similar color.
    % Mathematical formulation:
    \begin{equation}
        \mathcal{L}_{\text{affinity}} = \frac{1}{|\mathcal{N}|}
        \sum_{(i,j) \in \mathcal{N}}
        \mathbf{1}[|L_i - L_j| < \tau] \cdot \|ab_i - ab_j\|_1
    \end{equation}
    % where N is the set of adjacent pixel pairs and tau is the luminance threshold.

    \subsection{No Grey Loss (Novel)}
    % Penalizes the generator for leaving regions fully desaturated.
    % Mathematical formulation: TODO

\section{Training Procedure}
% Stage 1: pretraining generator with L1 loss only.
% Stage 2: full GAN adversarial training.
% Optimizer: Adam, lr_G = lr_D = 2e-4, beta1=0.5, beta2=0.999.
% Batch size, epochs, hardware (Tesla T4).


% ============================================================
\chapter{Experiments}
% ============================================================

\section{Experimental Setup}
% Hardware, dataset, training configuration.

\section{Evaluation Metrics}
% PSNR, SSIM, FID — define each and explain why it is relevant for colorization.

\section{Ablation Study}
% This is the most important section for the grade.
% A table comparing each loss combination:

\begin{table}[H]
    \centering
    \caption{Ablation study: effect of each loss component.}
    \begin{tabular}{lcccc}
        \toprule
        \textbf{Configuration} & \textbf{PSNR} $\uparrow$ & \textbf{SSIM} $\uparrow$ & \textbf{FID} $\downarrow$ \\
        \midrule
        Baseline (L1 + GAN)              & -- & -- & -- \\
        + TV Loss                        & -- & -- & -- \\
        + Contrast Loss                  & -- & -- & -- \\
        + Spatial Affinity Loss          & -- & -- & -- \\
        + No Grey Loss (full model)      & -- & -- & -- \\
        \bottomrule
    \end{tabular}
    \label{tab:ablation}
\end{table}

\section{Adversarial Robustness}
% FGSM and PGD experiments.
% Compare standard vs adversarially trained model under perturbation.

\section{Qualitative Results}
% Show the triplet figures: grayscale input | ground truth | colorized output.
% Pick representative examples: outdoor scenes, portraits, objects, etc.

\begin{figure}[H]
    \centering
    % \includegraphics[width=\textwidth]{../results/sample_outputs/example.png}
    \caption{Colorization results. Left: grayscale input. Center: ground truth. Right: model output.}
    \label{fig:results}
\end{figure}


% ============================================================
\chapter{Results and Discussion}
% ============================================================

\section{Quantitative Results}
\section{Effect of Novel Loss Functions}
\section{Limitations}
% Near-grayscale images in dataset.
% Dataset size (13k vs hundreds of thousands in SOTA).
% No perceptual/VGG loss.

\section{Future Work}


% ============================================================
\chapter{Conclusion}
% ============================================================

% Summarize: problem, approach, contributions, key results.
% One paragraph, no new information.


% ============================================================
\printbibliography[heading=bibintoc, title={References}]
% ============================================================


\end{document}
